% !TEX root = ../nc-quadric-surface.tex
\section{中心 standard 表示との比較}
\label{sec:standard-comparison}

\begin{proposition}[標識とシフトの移送]
\label{prop:standard-pointed}
$B=A/(g)$ を $\HH$ 対象の適合表示とし、$g$ は中心元、
$B$ は \cite{MN2025} の意味で standard とする。
$X$ を二つの ruling に対応する極大 Cohen--Macaulay 加群の一つとする。
すると \cite[Theorem 6.13]{MN2025} の三次元 AS 正則
$\mathbb Z$-algebra $D$（Gorenstein parameter 四）への同値 $F$ は
\[
 F(\pi B(i))\simeq\pi e_{-2i}D,\qquad
 F(\pi X(i))\simeq\pi e_{-2i-1}D
\]
を満たす。特に構造対象は $\pi e_0D$ に移る。
$D$ の二添字の周期移動から得る自己同値 $T$ に対して
$F(1)\simeq TF$ であり、$T^2\simeq S_D^{-1}[2]$ となる。
\end{proposition}

\begin{proof}
命題~\ref{prop:reconstruction} により $B$ は domain であり、
命題~\ref{prop:gldim-two} により $\qgr B$ の大域次元は有限である。
よって中心性と standard 性を加えれば
\cite[arXiv v1, Theorem 6.13]{MN2025} の仮定を満たす。
その構成は \cite[Theorem 5.16]{MU2021} の列
$E_{2i}=\pi B(i)$、$E_{2i+1}=\pi X(i)$ を用いる。
$D_{ij}=\Hom(E_{-j},E_{-i})$ の積は合成であり、
$j$ 成分が $\Hom(E_{-j},M)$ の右 $D$-加群を tails に送ると $F$ になる
\cite[arXiv v1, Definition 5.1, Theorems 5.8 and 6.4]{MN2025}。
従って $F(E_{-i})\simeq\pi e_iD$ である。

シフト $s=(1)$ は $s E_i\simeq E_{i+2}$ を与える。
これらの同型を使い、射に $s$ を作用させることで
$\alpha:D_{ij}\xrightarrow{\sim}D_{i-2,j-2}$ を得る。
右加群上で $(TN)_j=N_{j+2}$ とし、作用を $\alpha^{-1}$ で引き戻す。
写像
\[
 \Hom(E_{-j-2},M)\longrightarrow\Hom(E_{-j},sM),
 \qquad f\longmapsto s(f)
\]
（始域には上記の同一視を用いる）は合成と可換であり、
$TF\simeq Fs$ を与える。導来同値による Serre 関手の移送から
$T^2\simeq S_D^{-1}[2]$ が従う。

具体的に $p\le q$ に対して $B_{q-p}\to D_{2p,2q}$ は
$a\mapsto s^{-q}(a)$ であり、$p\le q\le r$ での合成は
$s^{-q}(a)\circ s^{-r}(b)=s^{-r}(s^{r-q}(a)\circ b)$ となる。
従って偶数添字部分は右規約の $\check B$ であり、
$R_n$ は $D_{0,4n}$ に対応する。積は周期移動と合成で与えられる。
\end{proof}

命題~\ref{prop:sklyanin-h2} の滑らかな Sklyanin 例では
\cite[Theorem 5.5]{MUW2026} が standard 性を保証するので、
命題~\ref{prop:standard-pointed} を適用できる。
両方の ruling を別々に用いて得た比較は、どちらも同じ $(\qgr B,\pi B)$ を回復する。
逆向きに任意の cubic $\mathbb Z$-algebra から出発した場合、
二添字の周期移動だけから $T^2\simeq S_D^{-1}[2]$ を結論することはできない。

% !TeX root = ../nc-quadric-surface-classification.tex

\subsection{三次正則表示を与える ruling}
\label{subsec:cubic-ruling}

\begin{proposition}[三次型の判定]
\label{prop:cubic-ruling}
命題~\ref{prop:standard-pointed} の中央標準表示 $B$ に対し、
二つの次数零生成 MCM 加群の一方 $X$ を選び、他方を $Y$ とする。
$E_{2n}=\pi B(n)$、$E_{2n+1}=\pi X(n)$ から作る
$D_{ij}=\Hom(E_{-j},E_{-i})$ は、次の積が同型である場合にちょうど
Van den Bergh の三次三次元 AS 正則 $\mathbb Z$-algebra となる。
\[
 \mu_X:\Hom(\pi X,\pi B(1))\otimes_k\Hom(\pi B,\pi X)
       \longrightarrow \Hom(\pi B,\pi B(1)),\qquad h\otimes x\longmapsto h\circ x .
\]
\end{proposition}

\begin{proof}
必要性は \cite[arXiv v3, \S 5.1, (Q2)--(Q3)]{VdB2011} の
距離二の積に関する記述から従う。三つの Hom はそれぞれ次元 $2,2,4$ である
\cite[Lemmas 5.6, 5.9]{MU2021}。

十分性を示す。$X(n)$ を標的とする射について、
\[
 0\longrightarrow Y(n-1)\longrightarrow B(n)^2\longrightarrow X(n)
   \longrightarrow0
\]
を使う \cite[Proposition 5.5(1)]{MU2021}。
tails で $E_j$（$j<2n+1$）から $\Hom$ を取ると、
$\Ext^1(E_j,\pi Y(n-1))=0$ なので、各射は
直前の $\pi B(n)$ を経由する射の和となる。消滅は
\cite[Lemmas 5.6(8), 5.9(2), 5.11]{MU2021} による。
$B$ の一次生成性と $\mu_X$ の全射性から、$\pi B(a)$ から $\pi B(b)$
（$b>a$）への射は、隣接する helix 対象を介した合成で張られる。
残る $\pi X(a)$ から $\pi B(b)$（$b>a$）への射に対しては、もう一方の列
$0\to X(-1)\to B^2\to Y\to0$ を $\Hom_B(-,B)$ で双対化する。
$Y$ は MCM なので $B^2\to\Hom_B(X,B)(1)$ は全射となる。
ここで次数付き内部 Hom を用いる。従って $X$ から $B$ への内部 Hom は、
次数一の二つの射で左 $B$-加群として生成され、その左作用は
$B$ のシフト間の射との合成である。tails との Hom の比較
\cite[Lemma 5.9(1)]{MU2021} を使えば、全ての射が隣接する helix 対象を
介した合成で張られる。

よって $D$ は隣接成分で生成され、各隣接成分は二次元である。
右単純加群の極小分解の初めは $P_{i+1}^2\to P_i$ であり、
左単純加群でも同じ隣接生成性を使える。
Mori--Nyman の AS 条件で大域次元 $3$、Gorenstein parameter $4$ なので、
$P_i=e_iD$ とすると右単純加群の最後の項は $P_{i+4}$ である。
双対複体が添字 $i+4$ の左単純加群の極小分解になることから、
残る項は $P_{i+3}^2$ である
\cite[arXiv v1, Theorem 4.5 and Remark 4.6]{MN2025}。
$D$ の成分次元は \cite[Lemmas 5.6, 5.9]{MU2021} の
Hilbert 級数から添字の差に関して二次で抑えられる。
これは \cite[arXiv v3, Definitions 4.1.1--4.1.2]{VdB2011} の三次正則型を与える。
\end{proof}

\begin{proposition}[Sklyanin pencil での判定]
\label{prop:sklyanin-cubic-ruling}
命題~\ref{prop:sklyanin-h2} の滑らかな対象に対し、
ruling のラベルを $w\in\{z,-z-2\tau\}$ と書く。
この ruling に属する線加群を $\operatorname{coker}(X(-1)\to B_z)$
と表示して得られる $X$ に対し、
\[
 \rank\mu_X=
 \begin{cases}
  4 & w\notin E[2],\\
  3 & w\in E[2].
 \end{cases}
\]
従って少なくとも一方の選択で、標識と平方根を保持した三次正則表示が得られる。
\end{proposition}

\begin{proof}
命題~\ref{prop:standard-pointed} の仮定は
\cite[Lemma 5.1 and Theorem 5.5]{MUW2026} と命題~\ref{prop:sklyanin-h2} により満たされる。
右加群での ruling と secant の対応には
\cite[Lemma 3.20, Proposition 3.21 and Lemmas 5.3--5.4]{MUW2026} を用いる。
ここでは補題~\ref{lem:sklyanin-opposite} の次数付き同型で
$B_z$ を通常の右加群表示 $T/(g_z)$ に移してから $w$ を読む。
固定した ruling は、$p+q=w$ を満たす secant $pq$ の線加群の族である。
実際、同じラベルの二つの線加群は Lemma 5.4 の列で共通の補線を取り、
異なるラベルの二つには同じ補題の列を直接適用できる。
標準性と Proposition 3.21 により、同じラベルと同じ ruling が対応する。

異なる一般の二つの secant $pq,p'q'$ を取り、四点は相異なるとする。
対応する線加群を定める $X(-1)\to B_z$ の二射は
二次元の Hom 空間の異なる一次元部分空間を与えるので、基底をなす。
それぞれの次数一像 $W,W'\subset (B_z)_1$ は secant 上で消える
線型形式の二次元空間であり、$\operatorname{im}\mu_X=W+W'$ である。
$W\cap W'\ne0$ は四点が同一平面上にあることと同値である。
楕円曲線の次数四の埋込みでは、四点が同一平面上にあることと
$p+q+p'+q'=0$ が同値である \cite[\S 10.1]{SVdB2013}。
従って共通部分が非零となるのは $2w=0$ の場合であり、
二つの secant は相異なるため、その場合の共通部分は一次元となる。
両ラベルが $E[2]$ に属すると $4\tau=0$ となり、
$\tau$ の無限位数に反する。
命題~\ref{prop:cubic-ruling} と命題~\ref{prop:standard-pointed} から結論を得る。
\end{proof}

$z=0$ は $\tau$ が無限位数のとき滑らかな例であり、
$w=0$ の選択では距離二の積が階数三、$w=-2\tau$ の選択では階数四である。
同じ構造対象と平方根から出発しても、どちらの ruling を取るかで
三次正則型の判定結果が変わる。


\subsection{制約を保つ Zhang twist}

\begin{lemma}[次数付き自己同型が tails 上で自明となる条件]
\label{lem:scalar-tails}
適合切断環 $B$ の次数付き自己同型 $\sigma$ に対し、
右作用の引き戻し $(-)_\sigma$ が $\qgr B$ 上で恒等と自然同型となるための
必要十分条件は、ある $\lambda\in k^\times$ に対して
$\sigma(b)=\lambda^d b$ ($b\in B_d$) となることである。
この次数スカラーを $\xi_\lambda$ と書く。
\end{lemma}
\begin{proof}
$\pi B(i)_\sigma\to\pi B(i)$ の基準同型は $b\mapsto\sigma^{-1}(b)$ である。
$\End(\pi B(i))=k$ より、恒等への任意の自然同型の $\pi B(i)$ 成分は
この基準同型の非零定数 $c_i$ 倍である。
$a\in B_1$ の左乗法 $\pi B(i)\to\pi B(i+1)$ に対する自然性から
$c_{i+1}\sigma^{-1}(a)=c_i a$ を得る。ここで左乗法は
$B_1\to\Hom(\pi B(i),\pi B(i+1))$ で単射である。
実際、非零 $a$ の左乗法は domain 性により非零の十分高い次数においても単射であり、
tails で零射にならない。一次生成により $\sigma$ は次数スカラーとなる。
逆に $\sigma=\xi_\lambda$ なら次数 $j$ の成分に $\lambda^{-j}$ を掛けることで
$M_\sigma\to M$ の自然同型が得られる。
\end{proof}

\begin{proposition}[固定表示での中心化と平方根性]
\label{prop:central-twist-obstruction}
適合表示 $B=A/(g)$ に対し、$ag=g\nu(a)$ で正規化自己同型 $\nu$ を定める。
$\sigma\in\Aut_{\mathrm{gr}}(A)$ が $\sigma(g)=c g$ を満たすとし、
積 $a\star b=a\sigma^{\deg a}(b)$ の Zhang twist を取る。
twist 後に $g$ が中心元となる条件は $\sigma^2=\xi_c\nu$ である。
一方 twist 後の商の次数シフトが反 Serre の平方根となる条件は
$\bar\sigma^2$ が $B$ の次数スカラーとなることである。
従って両条件が成り立つなら $\nu=\xi_1$ または $\nu=\xi_{-1}$ である。
\end{proposition}
\begin{proof}
$g$ の正則正規性が $\nu$ の存在と一意性を与える。
$a\in A_d$ に対し $a\star g=c^d g\nu(a)$、$g\star a=g\sigma^2(a)$
なので、中心性の判定は左から $g$ を消去すればよい。
これは $g$ の直線を保つ場合の
\cite[Remark 4.4]{MUW2026} の正規化公式である。

右加群の twist 同値 $F$ は作用を、次数 $j$ の $m$ に対して
$m\star b=m\bar\sigma^j(b)$ と変える。これは regular module を保つ。
$F(M)(1)$ の作用では $M(1)$ における次数 $j$ の元に元の次数 $j+1$ を使うので、
$F(M)(1)\simeq F((M(1))_{\bar\sigma})$ となる。
従って新しいシフトを $\qgr B$ に戻すと $s(-)_{\bar\sigma}$ であり、
その二乗は $s^2(-)_{\bar\sigma^2}$ となる。
$s^2\simeq S^{-1}[2]$ と Serre 関手の移送、および
補題~\ref{lem:scalar-tails} により平方根性の条件を得る。
両条件下で $\nu$ は $B$ 上で次数スカラー、従って
$A_1=B_1$ と一次生成性から $A$ 上でも次数スカラーである。
正規化の定義から $\nu(g)=g$ なので、その一次のスカラーは $\pm1$ である。
\end{proof}

\begin{example}
\label{ex:ore-central-twist}
命題~\ref{prop:ore-ambient-family} の $A_a$ では
$\nu(x)=a^{-1}x$、$\nu(y)=ay$、$\nu(u)=u$、$\nu(v)=v$ である。
$b^2=a$ を満たす $b$ を選べば
$(x,y,u,v)\mapsto (b^{-1}x,by,u,v)$ は関係と $g$ を保ち、
二乗が $\nu$ なので中心化できる。しかし $a\ne1$ なら $\nu$ は次数スカラーでなく、
この表示からはどの次数付き自己同型 twist でも
中心性とシフトの平方根性を同時には得られない。
同じ標識対象には $a=1$ の中心表示が存在するため、
中心表示の存在一般に対する反例は与えない。
\end{example}

\subsection{二添字周期が平方根となる条件}

\begin{proposition}
\label{prop:periodicity-gauge}
$D$ を命題~\ref{prop:standard-pointed} で構成した $\mathbb Z$-algebra とする。
より一般に \cite[Theorem 6.8]{MN2025} の coherent $\mathrm{ASF}^{++}$ 正則性の仮定を満たす
次元三・Gorenstein parameter 四の連結な $D$ で、
$D_{ij}\to\Hom(\pi e_jD,\pi e_iD)$ が同型である場合にも以下は成立する。
Nakayama 同型を $\nu:D_{ij}\to D_{i-4,j-4}$ とする。
二添字移動 $\alpha:D_{ij}\to D_{i-2,j-2}$ から得る自己同値 $T_\alpha$ が
$T_\alpha^2\simeq S_D^{-1}[2]$ を満たす必要十分条件は、
ある $c_i\in k^\times$ に対し
\[
 (\nu^{-1}\alpha^2)(a)=c_i c_j^{-1}a\qquad (a\in D_{ij})
\]
となることである。
また $\pi e_0D$ の反 Serre 環は次数 $n$ で $D_{-4n,0}$、
積は $a*b=\nu^n(a)b$ ($a\in D_{-4m,0}$、$b\in D_{-4n,0}$) となる。
\end{proposition}
\begin{proof}
命題~\ref{prop:standard-pointed} の $D$ に対する仮定は
\cite[Theorems 6.4, 6.13 and Proposition 4.11]{MN2025} による。
\cite[Theorem 6.8]{MN2025} の右規約では
$S_D^{-1}[2](N)=N(4)_{\nu^{-1}}$ であり、
$\pi e_iD$ を $\pi e_{i-4}D$ に移し、左乗法の射には $\nu$ を作用させる。
$T_\alpha^2$ は同じ添字移動を行い、射には $\alpha^2$ を作用させる。
二つの関手の比は添字を固定する自己同型 $\delta=\nu^{-1}\alpha^2$ による
作用の引き戻しである。
$\End(\pi e_iD)=k$ と標準 Hom 同型の下で
各左乗法に自然性を適用すると、$\delta$ が成分ごとに $c_i c_j^{-1}$ 倍となる。
逆にこの形なら、作用を $\delta^{-1}$ で引き戻した右加群から元の加群への同型は
次数 $j$ の成分で $c_j^{-1}$ を掛けて与えられる。従って tails 上でも自然同型を得る。

積は $S_D^{-1}[2]$ を $n$ 回左乗法に作用させてから合成するので、
表示された公式となる。次数 $n$ ごとに $\nu^{-n}$ で移すと
$D_{0,4n}$ 上で積は $a*b=a\nu^{-m}(b)$ となる。
特に適合表示が存在する場合、二次関係空間は
$D_{0,4}\otimes D_{0,4}\to D_{0,8}$、$a\otimes b\mapsto a\nu^{-1}(b)$
の核である。
\end{proof}
